\chapter{Введение: от набора секторов к одному рангоизменяющему полю}

\section{Почему новый том разрешено открыть}

Том VI завершился запретом продолжать прежнюю динамическую архитектуру по
инерции. Для открытия следующего тома требовался новый родительский кандидат,
который одновременно задаёт общий физический носитель, одно действие с одной
нормировкой и вычислимый тест запуска нулевого сектора. Новое название, новый
анзац дефекта или переназначение уже найденных чисел этого условия не
выполняют.

Том VII открывается с конкретным кандидатом. Его фундаментальным полем
является не заранее выбранный проектор $Q$, не фиксированный обменный мост
$K$ и не отдельное правило compacton-динамики. Основным объектом становится
физическая бимодульная одноформа переменного ранга
\begin{equation}
 \Phi(x)\in
 \Gamma\!\left(
 X,
 \operatorname{Hom}(\mathbb C^4,\operatorname{im}P_3)
 \widehat\otimes
 \mathcal Y_{\rm phys}
 \right),
 \qquad
 P_3=I_4-\frac14J_4,
 \label{eq:v7-fundamental-rank-field}
\end{equation}
где $X$ --- пространственно-временной носитель, а
$\mathcal Y_{\rm phys}=\Omega^1_{D_F}(\mathcal A_F)$ --- уже выделенный
физический модуль конечных одноформ. Поле $\Phi$ одновременно несёт
аффинный семейный, хиральный, пространственный и калибровочный смысл; его
Real-сопряжение включается тем же родителем.

Это отличие существенно. В Томе VI порядок $Q$ сначала выбирался как
отдельное состояние, после чего искалась связь с другими секторами. Здесь
проекторный порядок должен возникать как редуцированная квадратичная
характеристика самого поля $\Phi$. Поэтому ось, ранг и локальный профиль не
могут быть вставлены после вычисления.

\section{Наследуемое строгое ядро}

Том VII принимает без повышения статуса следующие результаты.
\begin{enumerate}
 \item Аффинное меню имеет каноническое разложение
 \begin{equation}
  \mathbb C^4=\operatorname{im}P_1\oplus\operatorname{im}P_3,
  \qquad
  P_1=\frac14J_4,
  \qquad
  \rank P_3=3.
 \end{equation}
 \item Физический конечный модуль одноформ существует, но прежний
 обменный мост класса пятнадцать не принадлежит ему как отдельная
 каноническая одноформа.
 \item Real-пара и её ориентированные классы $+15/-15$ классифицированы;
 класс пятнадцать является реестром полного коэффициентного пакета, а не
 числом частиц одного события.
 \item Поле порядка $Q\in\operatorname{Sym}_0(3,\mathbb R)$, его
 $\mathbb{RP}^2$-орбита и топологические дефекты построены как статический
 слой, но не получили общего динамического родителя.
 \item Унитарная симметричная динамика сохраняет характерные веса; чистая
 топология не создаёт вероятность рождения; прямая нить и compacton не
 являются готовыми конечными частицами.
 \item Любой переход между различными индексными секторами должен быть
 типизирован через спектральный поток, потерю ранга или иной явно заданный
 сингулярный слой.
\end{enumerate}

\section{Общий носитель}

Положим
\begin{equation}
 E_{\rm aff}
 :=\operatorname{Hom}(\mathbb C^4,\operatorname{im}P_3),
 \qquad
 \mathcal E_7
 :=S_X\widehat\otimes E_{\rm aff}
 \widehat\otimes\mathcal Y_{\rm phys},
 \label{eq:v7-common-carrier}
\end{equation}
где $S_X$ обозначает пространственный спинорный или решёточный переносный
носитель. Конфигурационное пространство состоит из допустимых связностей и
нечётных самосопряжённых Real-полей на $\mathcal E_7$.

Real-завершение фундаментального поля обозначается
\begin{equation}
 \widehat\Phi
 =\Phi+\varepsilon'J\Phi J^{-1},
 \qquad
 \widehat\Phi^*=\widehat\Phi.
 \label{eq:v7-real-completion}
\end{equation}
Нулевой вакуум равен $\Phi=0$. Ненулевые ранговые страты являются частями
одного линейного пространства полей, а не различными вручную склеенными
теориями.

При $\Phi\ne0$ семейное редуцированное состояние определяется формулой
\begin{equation}
 R_\Phi
 =\frac{\Tr_{\mathcal Y_{\rm phys}}(\Phi\Phi^*)}
 {\Tr(\Phi\Phi^*)},
 \qquad
 Q_\Phi=R_\Phi-\frac13I_3.
 \label{eq:v7-derived-order-parameter}
\end{equation}
В точке $\Phi=0$ нормированное состояние не определено, но исходное поле и
его действие определены. Поэтому вакуум не удаляется из конфигурационного
пространства нормировкой на заранее ненулевую амплитуду.

\section{Единая суперсвязность и действие}

Базовая суперсвязность имеет вид
\begin{equation}
 \mathbb A_\Phi
 =\mathbb A_0+\widehat\Phi,
 \qquad
 \mathbb A_0=\nabla_X\widehat\otimes1+1\widehat\otimes D_F,
 \label{eq:v7-parent-superconnection}
\end{equation}
а её кривизна равна
\begin{equation}
 \mathbb F_\Phi=\mathbb A_\Phi^2.
\end{equation}
Первый родительский кандидат Тома VII определяется одной нормой кривизны
\begin{equation}
 \mathcal S_7[\Phi,\nabla_X]
 =\int_X
 \operatorname{tr}_{\rm norm}
 \left(\mathbb F_\Phi^*\mathbb F_\Phi\right)d\mu_X.
 \label{eq:v7-parent-action}
\end{equation}
Здесь $\operatorname{tr}_{\rm norm}$ является одним нормированным следом по
всему конечному носителю. Независимые коэффициенты семейного, хирального,
Real- и радиационного секторов не вводятся. Общий положительный множитель
действия не влияет на классические уравнения и не используется как
наблюдаемая подгонка.

Действие \eqref{eq:v7-parent-action} является кандидатом, а не уже доказанной
физической теорией. Его право на дальнейшую разработку определяется первым
входным гейтом.

\section{Вычислимый тест запуска}

Для вариации $\eta$ положим
\begin{equation}
 L_0\eta=[\mathbb A_0,\widehat\eta]_{\rm s}.
\end{equation}
Тогда
\begin{equation}
 \mathbb F_{t\eta}
 =\mathbb F_0+tL_0\eta+t^2\widehat\eta^{2}.
\end{equation}
Условие стационарности нулевого сектора имеет вид
\begin{equation}
 \operatorname{Re}\langle\mathbb F_0,L_0\eta\rangle=0
 \quad
 \text{для всех физических }\eta,
 \label{eq:v7-vacuum-stationarity}
\end{equation}
а квадратичная форма физического гессиана равна
\begin{equation}
 \Hess_0\mathcal S_7(\eta,\eta)
 =2\|L_0\eta\|^2
 +4\operatorname{Re}
 \langle\mathbb F_0,\widehat\eta^{2}\rangle.
 \label{eq:v7-vacuum-hessian}
\end{equation}
Первый член неотрицателен. Единственный эндогенный источник отрицательной
моды --- кривизна уже заданного общего родителя, а не вставленная вручную
отрицательная масса.

\section{Рабочая и нулевая гипотезы}

\begin{description}
 \item[Рабочая гипотеза.] После полного gauge-, Real-, junk- и
 BRST/BV-факторирования оператор \eqref{eq:v7-vacuum-hessian} имеет
 физическую отрицательную моду. Её нелинейное развитие меняет ранг $\Phi$,
 проходит через сингулярную страту и насыщается конечным устойчивым
 объектом того же действия.

 \item[Нулевая гипотеза.] Допустимое пространство интертвинеров пусто или
 сводится к скалярной факторизации; либо единый след распадается на
 независимые веса; либо физический гессиан нулевого сектора неотрицателен.
 В любом из этих случаев кандидат закрывается до исследования частиц и
 наблюдаемых масштабов.
\end{description}

\section{Порядок работы Тома VII}

\begin{enumerate}
 \item Проверить типизацию общего носителя, Real-условие, калибровочную
 совместимость и наличие нескалярных интертвинеров.
 \item Зафиксировать единственный след и вывести полный физический гессиан
 нулевого сектора.
 \item Только при отрицательной моде построить нелинейную ветвь и проверить
 пересечение ранговой страты.
 \item Затем исследовать локализацию, устойчивость, открытый предел,
 скорость и масштаб.
 \item Лишь после этого разрешается обсуждать фермионное чтение, составные
 узлы и протонный сектор.
\end{enumerate}

\begin{nogocriterion}[Запрет обратного проектирования]
Поле $Q$, ось одноосности, ранг конечного объекта, compacton-фаза,
коэффициент $4\pi^2$, класс пятнадцать и дефицит $1/7$ не могут быть
подставлены в действие как цели. Они допустимы только как вычисленные
следствия одного поля $\Phi$.
\end{nogocriterion}

\begin{protocol}[Открытие Тома VII]
\begin{enumerate}
 \item новый фундаментальный кандидат: \textbf{поле $\Phi$ переменного
 ранга};
 \item общий носитель $P0$: \textbf{сформулирован};
 \item одно действие и одна нормировка $P1$: \textbf{сформулированы};
 \item вычислимый тест $P2$: \textbf{гессиан
 \eqref{eq:v7-vacuum-hessian}};
 \item физическая отрицательная мода: \textbf{не утверждается};
 \item устойчивый конечный объект: \textbf{не утверждается};
 \item наблюдаемая карта: \textbf{запрещена до динамического замыкания};
 \item следующий документ: \textbf{входной гейт допустимости
 рангоизменяющей суперсвязности}.
\end{enumerate}
\end{protocol}